\section{Finite local covers and cotangent tensors}
\label{loc:section}
\label{sec:local-tensor}

We first establish three geometric facts about the base germs: finite
smooth domination, rational smoothness, and the existence of maximal
finite quasi-\'etale charts.  The normalized source over such a chart
will also provide the direct-image and descent maps used later.
We then study reflexive cotangent tensors on one maximal chart through
a residual trace category.  Its tensor group acts on the fiber spaces
of that category.  The obstruction to a holomorphic connection
identifies the nontrivial irreducible summands of the cotangent fiber
representation as standard blocks for special linear and symplectic groups,
which enter the Euler calculation in
Section~\ref{sec:euler}.  A finite
map between normal spaces is called \emph{quasi-\'etale} if it is
\'etale outside a subset of codimension at least two in the target.
The geometric tensor argument uses complex analytic local rings of the
maximal charts; they are excellent and henselian, with residue field
$\C$.  The trace and split-lifting lemmas are stated for general normal
henselian local rings as well, so that they apply after completion.
Reflexive pullback means pullback followed by double dual.

\begin{lemma}[Finite smooth domination]
\label{loc:smooth-domination}
Every point of $B$ has an algebraic neighborhood, after an \'etale
base change if necessary, admitting a finite surjection from a smooth
variety of dimension $n$.  In particular, every analytic germ of $B$
is finitely dominated by a smooth germ of the same dimension.
Every such germ is analytically $\Q$-factorial.
\end{lemma}

\begin{proof}
Fix $b\in B$.  Choose sufficiently positive very ample divisors on the
smooth projective variety $X$.  By Bertini, a general intersection of
$n$ of them is smooth of dimension $n$.  They may simultaneously be
chosen so that their intersection with each component of $f^{-1}(b)$
is zero-dimensional and nonempty: the fiber has finitely many
components, all of dimension $n$, and successive general hyperplanes
cut each of them properly.  Denote the resulting smooth intersection
by $Y$.  The proper map $Y\to B$ is quasi-finite over $b$, and hence
finite over a neighborhood of $b$.  Its image contains that
neighborhood, by properness and the dimension theorem.  A local
branch at a point of $Y$ over $b$ gives the asserted finite map of
analytic germs.  If an individual algebraic branch is required, split
the finite algebra over the henselization and descend the corresponding
idempotents to an \'etale neighborhood.

Let $\pi:(Y,y)\to (B,b)$ be a finite dominating smooth germ, of degree
$m$, and let $D$ be a Weil divisor on $(B,b)$.  Pullback of Weil
divisors under a finite map is defined by the ramification indices of
the discrete valuations at prime divisors.  The divisor $\pi^*D$ is
principal on the smooth local germ, say $\operatorname{div}(g)$.
The valuation formula for the field norm gives
\[
 \operatorname{div}\bigl(\operatorname{Nm}_{\C(Y)/\C(B)}g\bigr)
       =\pi_*\pi^*D=mD.
\]
Thus $mD$ is Cartier.  The argument is analytic and therefore proves
the stronger local assertion needed later, rather than merely
algebraic $\Q$-factoriality.
\end{proof}

\begin{proposition}[Rational smoothness]
\label{loc:rational-smoothness}
The space $B$ is an oriented rational homology manifold of real
dimension $2n$.  Its rational intersection complex is
$\Q_B[n]$.  Consequently its ordinary rational cohomology has the
Betti numbers of $\PP^n$.  The local rational homology-manifold
assertion holds for every normal germ finitely dominated by a smooth
germ, in particular for the maximal charts constructed below.
\end{proposition}

\begin{proof}
We explain the trace at a branch point.  For a finite surjection
$\pi:Y\to U$ from a smooth $n$-fold to a normal connected
neighborhood, let $m_y$ be the local generic degree of the branch at
$y$.  A section $a$ of $\pi_*\Q_Y$ is locally constant on $Y$.
Set
\[
 \operatorname{Tr}_{\pi}(a)(b)
       =\sum_{y\in\pi^{-1}(b)}m_y a(y).
\]
This is locally constant.  Indeed, choose disjoint neighborhoods of
the finitely many points above $b$.  After shrinking downstairs, the
sum of the local generic degrees in each of these neighborhoods is
constant: it equals the degree over the common unbranched dense open
of the locally irreducible normal base.  The value of $a$ is constant
on each chosen neighborhood.  This proves the assertion and shows
that the trace is a morphism of sheaves.  Its composite with the
constant-section inclusion $\Q_U\to\pi_*\Q_Y$ is multiplication by
$\deg\pi$.  Therefore $\Q_U$ is a direct summand of
$\pi_*\Q_Y$.

Proper base change and the finiteness of the fibers give
$R^j\pi_*\Q_Y=0$ for $j>0$.  Applying the point costalk at $b$ to
the preceding direct summand, and using proper base change for
costalks, exhibits
\[
 H^k_{\{b\}}(U,\Q)
       \quad\text{as a direct summand of}\quad
 \bigoplus_{y\in\pi^{-1}(b)}H^k_{\{y\}}(Y,\Q).
\]
The latter group is zero unless $k=2n$.  The top local group on the
normal, hence locally irreducible, complex analytic germ $U$ is
one-dimensional: the complex orientation gives its fundamental
class, and the top fundamental cycles of a pure analytic germ are
indexed by its local irreducible components
\cite[Section~1.4, Theorem~3.2, Section~3.3 and Lemma~4.3]{BorelHaefliger1961}.
Here fundamental classes are taken in Borel--Moore homology of
representatives and passed to local groups. Equivalently, this
follows from an analytic triangulation and the oriented top strata.
Thus the costalk of $\Q_U$ is exactly $\Q[-2n]$ at every point,
with the compatible complex orientation.  This is the rational
homology-manifold assertion and identifies the intersection complex
with the shifted constant sheaf
\cite[Theorem~1.1(1)$\Leftrightarrow$(4)]{Massey2005}.
The final claim now follows from
the base-cohomology theorem recalled in
Proposition~\ref{prop:background}; the passage from intersection to
ordinary cohomology has been justified before using it.
\end{proof}

\begin{proposition}[Maximal quasi-\'etale charts and normalized source charts]
\label{loc:maximal-charts}
Every pointed germ $(B,b)$ has a connected finite normal Galois
quasi-\'etale cover $(C,c)$ which dominates every connected finite
quasi-\'etale cover of that germ.  The chart $(C,c)$ is klt, is
finitely dominated by a smooth germ, and has no nontrivial connected
finite quasi-\'etale cover.  The finite diagrams can be defined over
an algebraic \'etale neighborhood of $b$.

More generally, let $U\to B$ be an integral separated \'etale algebraic
neighborhood and let $\pi:C\to U$ be an integral finite normal
quasi-\'etale cover.  Put $X_U=X\times_B U$.  The base change
$X_U\times_U C$ has a unique irreducible component dominating $C$;
let $Z$ be its normalization and $h:Z\to C$ its structure map.
Then $h$ is projective and equidimensional of relative dimension $n$,
has connected fibers, and satisfies $h_*\cO_Z=\cO_C$.  The finite map
$p:Z\to X_U$ is \'etale.  Consequently $Z$ is smooth, the pullback of
the symplectic form on $X$ is symplectic on $Z$, its canonical bundle is
trivial, and the fibers of $h$ are Lagrangian.  If $\pi$ is Galois with
group $\Gamma$, then $p$ is a $\Gamma$-torsor.  Ordinary \'etale base
change on $C$ preserves projectivity, equidimensionality, connected
fibers, the equality $h_*\cO_Z=\cO_C$, and the smooth symplectic source.
The torsor in the statement belongs to the original finite chart.
\end{proposition}

\begin{proof}
Work first over the strict henselization of the algebraic germ and
choose a finite smooth slice as in
Lemma~\ref{loc:smooth-domination}.  Choose a connected local factor
of its finite algebra, with fraction field $L$ over the base field
$K$.  If $K'/K$ defines a finite quasi-\'etale normal cover, the
normalization of its pullback to this slice is unramified outside a
set of codimension at least two.  Finiteness preserves that
codimension bound.  The source of the normalized map is normal and
the target is regular, so purity of the branch locus
\cite[Tag 0BMB]{StacksProject} makes the map \'etale everywhere.
A connected finite \'etale cover of a strictly henselian regular
local scheme is trivial.  A component of the split pullback
therefore embeds $K'$ into $L$ over $K$.

Let $M/K$ be a finite Galois closure of $L/K$, let $W$ be the
normalization of the base in $M$, and write $\Gamma$ for its Galois
group.  Normalizations are finite by excellence.  Let $H\lhd\Gamma$
be the normal subgroup generated by the inertia groups of all prime
divisors of $W$.  The finite quotient $C=W/H$ is unramified over
every prime divisor of the normal base and hence is quasi-\'etale.
Conversely, an intermediate extension which is quasi-\'etale has
trivial divisorial inertia; its defining subgroup contains every
conjugate of each inertia group, and hence contains $H$.  Every
quasi-\'etale cover is consequently dominated by $C$.  A further
quasi-\'etale cover of $C$ would give one of the base and is therefore
trivial.  Applying the same split-pullback argument to $C$ shows
that it too is finitely dominated by the smooth slice.

All finite algebras, group actions, quotient maps and their
identities descend from the strict henselization to an \'etale
neighborhood.  Shrinking removes branch divisors not meeting the
marked fiber.  This construction also gives the analytic maximum.
Indeed, the normalized pullback of any finite analytic
quasi-\'etale cover to the smooth germ is \'etale by analytic purity
and splits because a smooth germ is simply connected.  One may
also obtain this directly: a small smooth ball minus an analytic
set of complex codimension at least two is simply connected, by
general position for disks against a stratification; its finite
cover splits there, and normalization extends that splitting
uniquely.  The analytic cover is thus
an intermediate cover of the analytification of the same finite
Galois closure.  Finite normalizations and finite-group quotients
commute here with analytification; the local normal components and
their degrees agree.  The divisorial inertia groups through the
marked point are therefore the same in both constructions.  Finally
$C$ is klt.  Indeed, writing $\pi:C\to B$, we have
$K_C=\pi^*K_B$ in codimension one, and for a
divisorial valuation $E$ over $C$ restricting with ramification index
$r$ to a divisorial valuation $F$ over $B$, the canonical ramification
formula on normalized birational models gives
$A_C(E)=rA_B(F)$.  The latter is positive because $B$ is klt.

We prove the source assertions, including connectedness.  The connected
fibers of $f$ and normality of $B$ give $f_*\cO_X=\cO_B$ by Stein
factorization.  This equality persists after the flat base change to
$U$.  Let $K=\C(U)$, $L=\C(C)$, and let $F$ be the generic fiber of
$X_U\to U$.  Generic smoothness makes $F$ smooth.  Proper flat field
base change gives
\[
 H^0(F_{\overline K},\cO_{F_{\overline K}})=\overline K.
\]
Thus $F_{\overline K}$ is connected.  A connected smooth scheme over
an algebraically closed field is integral, since its irreducible
components are disjoint.  In this setting it is an abelian variety
after an origin is chosen, by Proposition~\ref{prop:background}.
In particular $F$ is geometrically integral, so $F_L$ is integral
and normal and has $H^0(F_L,\cO_{F_L})=L$.  This proves uniqueness of
the component dominating $C$ and identifies the generic fiber of its
normalization $Z$ with $F_L$.

The normalization is finite by excellence, so $h$ is projective.
Its Stein algebra $h_*\cO_Z$ is finite over $\cO_C$ and embeds in the
generic algebra $L$: a section on the integral space $Z$ is determined
by its restriction to the generic fiber.  Every such section is
integral over $\cO_C$; normality of $C$ therefore gives
$h_*\cO_Z=\cO_C$.  Stein factorization now gives connected fibers.
This part of the argument used only that $C\to U$ is finite and
dominant, with $C$ normal and integral.

Let $D\subset U$ have codimension at least two and contain the branch
locus of $\pi$.  Equidimensionality of $f$ gives
\[
 \dim f_U^{-1}(D)\leq\dim D+n\leq2n-2.
\]
The map $p$ is finite over the smooth space $X_U$ and is \'etale away
from $f_U^{-1}(D)$.  Purity makes it \'etale everywhere, so $Z$ is
smooth and carries the pulled-back symplectic form.  Each component of
a fiber of $h$ has dimension at most $n$, because its finite image lies
in a fiber of $f_U$.  The fiber dimension theorem gives the reverse
inequality.  Hence $h$ is equidimensional of relative dimension $n$;
the same finite images show that its fiber components are Lagrangian.
The symplectic form trivializes the canonical bundle of $Z$.

If $\pi$ is Galois, its deck action lifts to $Z$.  Over the dense
unbranched locus the map $p$ is a deck torsor.  Both sides of the torsor
map $\Gamma\times Z\to Z\times_{X_U}Z$ are finite \'etale over $Z$,
so their isomorphism on that dense locus extends over $Z$.  This proves
freeness of the action also over special fibers.  Subsequent \'etale
refinements of the object chart use ordinary base change of $h$ and
preserve its projectivity, equidimensionality, connected fibers and
smooth symplectic source.  The torsor descent is retained from the
finite chart before that refinement.
\end{proof}

\subsection{The trace quotient}

Fix a maximal chart and let $(R,\mathfrak m)$ be its analytic local
ring.  Put $V=\Omega_R^{[1]}$.  On reflexive modules we use
$M\mathbin{\widehat\otimes}N=(M\otimes_RN)^{**}$ and
$M^\vee=\Hom_R(M,R)$.  These are a symmetric rigid tensor product
and dual: on a common open set containing every point of codimension
one they are the ordinary bundle operations, and all identities
extend uniquely by reflexivity.  In particular the categorical trace
is the usual generic trace, which lies in $R$ by normality.  We use
the full additive, idempotent-complete tensor subcategory $\mathcal C$
generated by $V$ and $V^\vee$.  Every object in this subcategory is a vector
bundle on the smooth locus of the chart.

There are two uses of endomorphisms below.  The ordinary finite
$R$-algebra $A_M=\End_R(M)$ consists of morphisms of the module $M$.
The internal endomorphism object is
\[
 \underline{\End}(M)=M\widehat\otimes M^\vee.
\]
Its underlying reflexive module is canonically $\Hom_R(M,M)$: this is
the usual bundle identification in codimension one, extended
reflexively.  We use the underline when this module is a coefficient
object in $\mathcal C$, rather than the algebra of morphisms of $M$.

\begin{lemma}[The trace ideal is the Jacobson radical]
\label{loc:trace-radical}
Let $(R,\mathfrak m)$ be any noetherian normal henselian local
$\C$-algebra with residue field $\C$.  For finite reflexive modules
$M,N$, define
\[
 \mathcal R(M,N)=
 \{a\in\Hom_R(M,N):\tr(ba)\in\mathfrak m
                 \text{ for every }b\in\Hom_R(N,M)\}.
\]
These spaces form a tensor ideal.  For every $M$,
$\mathcal R(M,M)$ is the Jacobson radical of the finite
$R$-algebra $\End_R(M)$.  The quotient category is semisimple
abelian over $\C$, has finite-dimensional Hom spaces, and its
categorical dimensions are the positive integer ranks of its
nonzero objects.
\end{lemma}

\begin{proof}
Cyclicity of trace proves stability under composition.  Stability
under tensor product follows by taking a partial trace: for a map
$a:M\to N$ and a test map
$b:N\widehat\otimes P\to M\widehat\otimes P$, the trace of
$b(a\widehat\otimes\id_P)$ is the trace of $a$ composed with the
partial trace of $b$.  These formulas hold for bundles in codimension
one and therefore over $R$.  Duality is compatible with the same
ideal.

Write $A=A_M=\End_R(M)$ and $J=\rad A$.  The algebra $A$ is finite over
the henselian local ring $R$, so $\mathfrak m A\subset J$,
$A/J$ is a finite-dimensional semisimple $\C$-algebra, and
$J^q\subset\mathfrak m A$ for some $q$.  We recall why idempotents
lift in this setting.  First lift an idempotent through the nilpotent
radical of $A/\mathfrak m A$.  For an ensuing lift $a\in A$ of an
idempotent modulo $\mathfrak m A$, the commutative finite algebra
$R[a]$ splits into the factors supported at $0$ and $1$ in its
closed fiber by henselianity.  The resulting idempotent lifts the
given one.  Induction in complementary corners lifts an entire
orthogonal system.  Thus primitive idempotents of $A/J$ lift to
actual direct summands of $M$.

For $x\in J$, all sufficiently large powers $x^j$ belong to
$\mathfrak m A$, so their traces vanish modulo $\mathfrak m$.
The generic characteristic polynomial of $x$ has coefficients in
$R$: this is true at every height-one localization, where $M$ is
free, and $R$ is the intersection of those valuation rings.  If
$\lambda_1,\ldots,\lambda_r$ are the roots of its reduction over
$\C$, Newton identities give
$\sum_i\lambda_i^j=0$ for every sufficiently large $j$.
Grouping the distinct nonzero roots, a Vandermonde matrix applied to
consecutive such powers forces all their multiplicities to be zero.
There are consequently no nonzero roots, and $\tr(x)=0$ modulo
$\mathfrak m$.  The same holds for $yx$ for any $y\in A$, because
$J$ is an ideal.  Hence $J\subset\mathcal R(M,M)$.

The residual trace now descends to $A/J$.  On each simple block
$\operatorname{Mat}_t(\C)$ it is a scalar multiple of the ordinary
matrix trace, by cyclicity.  Evaluate it on a primitive diagonal
idempotent.  Its lift projects onto a nonzero reflexive direct
summand $M_0$, so the scalar is $\rk(M_0)>0$.  The trace pairing is
therefore nondegenerate on every block.  This proves the reverse
inclusion and the equality with $J$.

Apply the argument to finite direct sums of objects.  Their
endomorphism algebras in the quotient are semisimple matrix algebras,
and their primitive projectors lift.  In particular every object is
a finite sum of simple objects, distinct simple objects have no
morphisms, and the endomorphisms of each simple object are $\C$.
This also proves that the quotient is abelian and that all its exact
sequences split.  The ideal contains $\mathfrak m\Hom_R(M,N)$,
so its Hom spaces are finite-dimensional.  Finally the trace of the
identity is $\rk M$, as asserted.
\end{proof}

\begin{lemma}[Lifting split maps]
\label{loc:split-lift}
Over any ring in Lemma~\ref{loc:trace-radical},
if a morphism of reflexive modules becomes a split injection in the
trace quotient, it is already a split injection of reflexive modules.
An isomorphism in the quotient lifts to an isomorphism.  In
particular, an object is a sum of copies of the tensor unit in the
quotient if and only if its underlying module is free.
\end{lemma}

\begin{proof}
Let $a:M\to N$ and lift a left inverse in the quotient to
$b:N\to M$.  Then $ba=\id_M+r$ with $r\in\rad\End_R(M)$,
so $ba$ is invertible.  Replacing $b$ by $(ba)^{-1}b$ gives an
actual left inverse.  If the quotient map is an isomorphism, apply
the same argument to both composites of lifts of it and its inverse;
the lift is an isomorphism.  Apply this to an isomorphism with $R^r$
for the final statement.
\end{proof}

Return now to the analytic maximal chart and the category $\mathcal C$
fixed above.

\begin{proposition}[The residual tensor group]
\label{loc:tannakian}
The residual category $\overline{\mathcal C}=\mathcal C/\mathcal R$
admits a fiber functor $\omega$ to complex vector spaces.  For a choice
of $\omega$, tensor reconstruction identifies it with
$\operatorname{Rep}_{\C}(G)$, where
$G=\Aut^\otimes(\omega)$ is a reductive algebraic subgroup of
$\operatorname{GL}(V_{\mathrm{fib}})$ and
$V_{\mathrm{fib}}=\omega(\overline V)$.  This representation is
faithful and has dimension $n$.
\end{proposition}

\begin{proof}
Lemma~\ref{loc:trace-radical} gives an abelian rigid symmetric
$\C$-linear category with $\End({\bf1})=\C$.  Tensor product is
exact because every exact sequence splits.  For an object of rank
$r$, its $(r+1)$st exterior power is zero: it is zero on the common
bundle locus and hence its reflexive extension is zero.  Deligne's
criterion \cite[Theorem 7.1]{Deligne1990} therefore applies.  The
category is tensor-generated by one object; Deligne's
Corollary~6.20 supplies a fiber functor over a finite extension of
$\C$, hence over $\C$ itself.  Tensor reconstruction gives the
stated group.  A tensor generator makes it an algebraic subgroup of
its general linear group, and semisimplicity of the representation
category makes it reductive in characteristic zero.  The fiber
dimension equals the categorical dimension, namely $\rk V=n$.
\end{proof}

We call $G$ the residual tensor group of the chosen fiber functor.
Its representations are the fiber spaces of $\overline{\mathcal C}$.
Finite deck groups, denoted by $\Gamma$, act on the geometric charts.
For example, if $M_{\mathrm{fib}}=\omega(\overline M)$, then
\[
 \omega\bigl(\overline{\underline{\End}(M)}\bigr)
       =\End_{\C}(M_{\mathrm{fib}}),\qquad
 \End_{\overline{\mathcal C}}(\overline M)
       =\End_G(M_{\mathrm{fib}}).
\]
The first space is the full internal coefficient representation, with
the conjugation action; the second is its space of invariant vectors.

\subsection{Connections before and after the trace quotient}

For an object $M$ let $M_U$ be its vector bundle on the smooth locus
$U$ of a sufficiently small representative of the chart.  Its
ordinary Atiyah class \cite{Atiyah1957} is
\[
 a(M_U)\in H^1\bigl(U,\Omega_U^1\otimes\End(M_U)\bigr).
\]
It vanishes precisely when $M_U$ has a holomorphic connection;
integrability is not part of this assertion.  We take the direct
limit over shrinking representatives when discussing this class on
a germ.  All cohomology in the following lemmas is this ordinary
sheaf cohomology.

\begin{lemma}[A connection gives freeness on a maximal chart]
\label{loc:connection-free}
If $a(M_U)=0$, then $M$ is a free $R$-module.
\end{lemma}

\begin{proof}
Pull a connection on $M_U$ back to a finite smooth dominating germ
$\pi:(Y,y)\to(C,c)$.  The complement of $\pi^{-1}(U)$ has
codimension at least two.  On $Y$ let
$F=(\pi^*M)^{**}$.  In local smooth coordinates each coefficient
of the connection extends across that complement by the Hartogs
property of $F$.  Its target $F\otimes\Omega_Y^1$ is reflexive
because $\Omega_Y^1$ is locally free.  Thus the connection extends
to $F$ and still satisfies Leibniz's rule.

Here is the elementary freeness argument, which does not require a
flat connection.  For a coordinate derivation $\partial$, choose a
presentation of $F$ and lift the operator $\nabla_{\partial}$ to
the free module of generators.  If $A$ is the corresponding matrix
on generators and $P$ a matrix of relations, stability of the
relations gives a matrix $Q$ with
\[
                 \partial P+AP=PQ.
\]
The derivative of every minor of $P$ is therefore a linear
combination of minors of the same size.  Each Fitting ideal of $F$
is stable under every coordinate derivation.  A nonzero ideal in a
smooth characteristic-zero local ring with this property is the
unit ideal: repeatedly differentiating a nonzero element of minimal
order produces an element with nonzero constant term.  If $r$ is
the generic rank of $F$, its $r$th Fitting ideal is nonzero and hence
is the unit ideal.  Thus $F$ is generated by $r$ elements.  The
resulting surjection from a free module of rank $r$ has a torsion
kernel, which is zero because it is a submodule of a free module
over a domain.  This proves that $F$ is free.

Take a normal Galois closure $W\to C$ of $Y\to C$, with group
$\Gamma$.  Its normal domain is local over the henselian germ,
so the whole group fixes the unique marked point.  The
reflexive pullback $P$ of $M$ to $W$ is free: it agrees in
codimension one with the pullback of the free module $F$, and both
reflexive extensions agree.  It has its natural semilinear
$\Gamma$-action.  At a prime divisor, divisorial inertia acts
trivially on its restriction to that divisor, because there $P$
is the ordinary pullback of a bundle from $C$.  The same inertia
therefore acts trivially on the fiber of $P$ at the marked point.
To see the last specialization explicitly, for an inertia element
of order $m$, restrict $P$ to its fixed locus and apply the
idempotents
\[
       \frac1m\sum_{j=0}^{m-1}\zeta^{-j}g^j.
\]
Their ranks are locally constant.  On the divisor all nonidentity
eigenvalues have rank zero, hence they have rank zero at its
specialization to the marked point as well.

Since $C$ has no further quasi-\'etale cover, these inertia groups
normally generate $\Gamma$, by the construction in
Proposition~\ref{loc:maximal-charts}.  The action on the marked
fiber of $P$ is consequently trivial.  Lift a basis of that fiber
to sections of $P$ and average over $\Gamma$.  The averaged
sections still lift the basis and hence, by Nakayama, form a basis
of the free module $P$.  This basis is $\Gamma$-invariant.
The invariant basis identifies $P^\Gamma$ with $R^r$, because the
invariants of the normal cover ring are $R$.  Thus $P^\Gamma$ is
reflexive.  The natural map $M\to P^\Gamma$ is an isomorphism in
codimension one by ordinary finite descent for the pullback of a
bundle.  Both modules are reflexive, so this isomorphism extends over
the normal germ.  Hence $M=P^\Gamma$ is free, as claimed.
\end{proof}

\begin{lemma}[Tensor tests for the ordinary Atiyah class]
\label{loc:atiyah-tests}
Let $M$ be an object, and let $t_j$ be actual sections of finitely
many functorial mixed tensor powers $T_j(M)$ of $M$ and $M^\vee$ on
the chart.  Direct sums and the usual symmetric or alternating
projectors may be used; each construction carries the natural action
of every endomorphism of $M$.  Put $Q=\bigoplus_jT_j(M)$ and define
\[
 \Phi_M:\underline{\End}(M)\longrightarrow Q,
       \qquad A\longmapsto(A\cdot t_j)_j
\]
by the infinitesimal tensor action.  The coefficient map for the
Atiyah class is the morphism in $\mathcal C$
\begin{equation}\label{loc:psi-coefficient}
 \Psi_M=\id_V\widehat\otimes\Phi_M:
 V\widehat\otimes\underline{\End}(M)
       \longrightarrow V\widehat\otimes Q.
\end{equation}
On the smooth locus its ordinary bundle restriction satisfies
$(\Psi_{M,U})_*a(M_U)=0$ in $H^1(U,(V\widehat\otimes Q)_U)$.
If $\overline{\Psi_M}$ is injective in $\overline{\mathcal C}$, then
$a(M_U)=0$.
\end{lemma}

\begin{proof}
The natural tensor action defining $\Phi_M$ is a bundle map on $U$;
its reflexive extension is the displayed morphism in $\mathcal C$.
Choose local connections $\nabla_\alpha$ on $M_U$.
Their differences represent $a(M_U)$.  On $T_j(M_U)$ the
corresponding differences are the infinitesimal tensor actions of
$\nabla_\alpha-\nabla_\beta$.  Evaluating on the global section
$t_j$ gives the coboundary of the local sections
$\nabla_\alpha t_j$.  This proves the asserted vanishing after
restricting \eqref{loc:psi-coefficient} to bundles on $U$ and taking
ordinary $H^1$.

An injection in the semisimple quotient has a left inverse there.
Lemma~\ref{loc:split-lift} lifts it to an actual left inverse
between reflexive coefficient modules.  Restriction to $U$ gives
a split map of bundles, which induces an injection on ordinary
$H^1$.  Its kernel contains $a(M_U)$, so that class is zero.
The argument uses only that the individual $t_j$ are actual
sections.  No algebraic identities among their images in the
quotient need hold for the chosen lifts.
\end{proof}

\begin{proposition}[Connectedness and semisimplicity]
\label{loc:connected-semisimple}
The reductive group $G$ of Proposition~\ref{loc:tannakian} is
connected and semisimple.
\end{proposition}

\begin{proof}
For a reductive subgroup of a general linear group, the invariant
mixed tensors have that subgroup as their simultaneous pointwise
stabilizer.  Indeed, projectors onto direct summands and all maps
between tensor constructions are themselves such tensors.
Preserving them is precisely the condition to give a tensor
automorphism of the forgetful functor, and tensor reconstruction
identifies those automorphisms with the group.  Differentiating,
its Lie algebra is the common kernel of the infinitesimal
tensor-action maps.  Since the endomorphism space is
finite-dimensional, finitely many tensors suffice for this last
kernel.

Suppose a representation $N$ of $G$ has finite image.  Realize it
as a direct summand of tensor constructions in
$V_{\mathrm{fib}}$, and lift its projector to obtain an actual
module $M$.  Choose finitely many invariant mixed tensors of $N$
whose infinitesimal stabilizer is the zero Lie algebra.  By the
definition of the quotient Hom spaces these tensors lift to actual
maps from $R$, hence to actual sections on the chart.  The map
$\omega(\overline{\Phi_M})$ of Lemma~\ref{loc:atiyah-tests} has as
its kernel this infinitesimal stabilizer inside the full space
$\End_{\C}(N)$, and is therefore injective.  Thus
$\overline{\Phi_M}$ is injective.  Tensoring it with $\overline V$
preserves injectivity, so $\overline{\Psi_M}$ is injective and that
lemma and Lemma~\ref{loc:connection-free} make $M$ free.
Therefore $N$ is trivial.  The regular representation of the
finite component group $G/G^\circ$ is one such representation.
It is trivial only if that component group is trivial.  Hence $G$
is connected.

A character of $G$ corresponds to a rank-one reflexive module on
the chart.  By analytic $\Q$-factoriality some positive reflexive
tensor power of that module is free.  Thus the character has
finite order.  A connected algebraic group has no nontrivial
finite-image character.  The character group of $G$ is therefore
zero.  For a connected reductive group its quotient by the derived
group is a torus, whose characters detect it; that torus must be
trivial.  Thus $G$ is semisimple.
\end{proof}

\subsection{The prolongation forced by a nontrivial cotangent block}

For a vector space $T$ and a Lie subalgebra
$\mathfrak h\subset\End(T)$, set
\[
 \mathfrak h^{(1)}=
 \{A\in\Sym^2T^*\otimes T:
                      A(x,{-})\in\mathfrak h\text{ for every }x\in T\}.
\]
For a reflexive module $M$, let $\Sym^{[2]}M$ denote the image of the symmetric idempotent
$(1+\tau)/2$ on $M\widehat\otimes M$.  It is a reflexive direct
summand and restricts to the usual symmetric square on the bundle
locus.  For a block $V_i$ below with fiber
$L_i=\omega(\overline V_i)$ and $T_i=L_i^\vee$, the fiber of its
coefficient module is
\[
 \omega\bigl(\overline{\Sym^{[2]}V_i\widehat\otimes V_i^\vee}\bigr)
       =\Sym^2L_i\otimes L_i^\vee
       =\Sym^2T_i^*\otimes T_i.
\]
The intersection defining a prolongation will be taken in this fiber
space.
We use Lie's triangularization theorem and complete reducibility
for finite-dimensional complex representations, in the forms
\cite[Theorem 2.2(2), Corollary 2.2(2), Theorem 2.9]{Cap2016}.
The prolongation spaces belong to the Singer--Sternberg framework
\cite{SingerSternberg1965}. The infinite-type rank-one criterion is due
to Guillemin--Quillen--Sternberg \cite{GuilleminQuillenSternberg1967};
see also the proof and attribution in
\cite[Section~2.2 and Theorem~2.3]{OttazziWarhurst2011}.
We give the prolongation argument in full, including the separate
semisimple-irreducible step from nonzero first prolongation to infinite
type.

\begin{lemma}[One-block connection obstruction]
\label{loc:block-obstruction}
Lift a decomposition of $V_{\mathrm{fib}}$ into irreducible
$G$-modules to a decomposition $V=\bigoplus_iV_i$ of reflexive
modules.  Let
\[
 \mathfrak h_i=
    \mathfrak g\cap\End(V_{i,\mathrm{fib}})
       \subset\End(V_{\mathrm{fib}}),
       \qquad\mathfrak g=\operatorname{Lie}(G),
\]
where the endomorphisms in this intersection are zero on every
other block.  Regard $\mathfrak h_i$ on
$T_i=V_{i,\mathrm{fib}}^\vee$ by the dual action.  If
$\mathfrak h_i^{(1)}=0$, then $V_i$ is free.  Consequently every
nontrivial irreducible block has
$\mathfrak h_i^{(1)}\ne0$.
\end{lemma}

\begin{proof}
Local coordinate connections on $\Omega_U^1$ are torsion-free.
Their differences show that $a(V_U)$ belongs to the direct summand
\[
 H^1\bigl(U,\Sym^2V_U\otimes V_U^\vee\bigr)
       \subset H^1\bigl(U,V_U\otimes\End(V_U)\bigr).
\]
On the other hand, choose local connections preserving the actual
splitting $V_U=\bigoplus_i(V_i)_U$.  The off-diagonal
endomorphism components of the same cohomology class are zero.
These two conditions imply
\begin{equation}
 a(V_U)=\sum_i a_i,\qquad
 a_i\in H^1\bigl(U,\Sym^2(V_i)_U\otimes(V_i)_U^\vee\bigr).
 \label{loc:diagonal-atiyah}
\end{equation}
Here is the index check.  Block preservation allows only
coefficients $V_a\otimes V_b\otimes V_b^\vee$.  Symmetry
interchanges the first two covariant slots.  If $a\ne b$, its
image is an off-diagonal coefficient
$V_b\otimes V_a\otimes V_b^\vee$, whose class is zero.
Because these are actual direct sums of bundles, the original
mixed component is also zero.  This leaves precisely the terms in
\eqref{loc:diagonal-atiyah}.  In particular $a_i$ is the class
$a((V_i)_U)$ under the indicated inclusion of its coefficients.

Choose lifts $t_j$ of finitely many $G$-invariant tensors in natural
mixed powers of $V_{\mathrm{fib}}$ cutting out $\mathfrak g$
infinitesimally.  Use them to define $Q$ and $\Phi_V$ as in
Lemma~\ref{loc:atiyah-tests}.  For the $i$th block set
\[
 S_i=\Sym^{[2]}V_i\widehat\otimes V_i^\vee.
\]
The symmetric inclusion in the first two slots, followed by the
diagonal inclusion of the $i$th endomorphism block, gives
\[
 \jmath_i:S_i\longrightarrow
          V_i\widehat\otimes\underline{\End}(V).
\]
The required one-block coefficient map is
\begin{equation}\label{loc:kappa-coefficient}
 \kappa_i=(\id_{V_i}\widehat\otimes\Phi_V)\circ\jmath_i:
 S_i\longrightarrow V_i\widehat\otimes Q.
\end{equation}
On $U$ the source of this map is
$\Sym^2(V_i)_U\otimes(V_i)_U^\vee$.  Projecting the first slot of
$(\Psi_{V,U})_*a(V_U)=0$ onto $(V_i)_U$ and using
\eqref{loc:diagonal-atiyah} gives
$(\kappa_{i,U})_*a_i=0$ in ordinary $H^1$.

Write $L=V_{\mathrm{fib}}$ and $L_i=V_{i,\mathrm{fib}}$.
The kernel of $\omega(\overline{\Phi_V})$ is $\mathfrak g$ inside
$\End_{\C}(L)$.  Hence the fiber kernel of
\eqref{loc:kappa-coefficient} is the explicit intersection
\begin{equation}\label{loc:kappa-kernel}
 \ker\omega(\overline{\kappa_i})=
 \bigl(\Sym^2L_i\otimes L_i^\vee\bigr)
          \cap\bigl(L_i\otimes\mathfrak g\bigr)
       \quad\text{inside }L_i\otimes\End_{\C}(L).
\end{equation}
Partial evaluation of the first covariant slot by $x\in T_i=L_i^\vee$
is supported on the $i$th endomorphism block.  It belongs to
$\mathfrak g$ precisely when it belongs to $\mathfrak h_i$.
The symmetry of the first two slots is the defining symmetry in
$\mathfrak h_i^{(1)}$ on $T_i$.  Passing from the action on $L_i$
to the dual action on $T_i$ transposes and changes sign, which leaves
this linear kernel unchanged.  Thus \eqref{loc:kappa-kernel} identifies
with $\mathfrak h_i^{(1)}$.

If the kernel is zero, faithfulness of $\omega$ makes
$\overline{\kappa_i}$ injective, hence split injective in the semisimple
quotient.  Lemma~\ref{loc:split-lift} lifts a left inverse for
$\kappa_i$ between the reflexive coefficient modules.  Restricting
that left inverse to $U$ makes the induced ordinary $H^1$ map injective.
Thus $a_i=0$ in ordinary $H^1$, and
Lemma~\ref{loc:connection-free} makes $V_i$ free.  A free
irreducible object is the tensor unit.  This proves the assertion.
\end{proof}

\begin{lemma}[An auxiliary tensor factor kills prolongation]
\label{loc:auxiliary-factor}
Suppose $T=U\otimes W$, $\dim W>1$, and
$\mathfrak h\subset\End(U)\otimes\id_W$.  Then
$\mathfrak h^{(1)}=0$.
\end{lemma}

\begin{proof}
Write a prolongation tensor as $A(x,y)=A_x(y)$ with
$A_x\in\mathfrak h$.  Fix $x=u\otimes w_1$ with $w_1\ne0$
and choose $w_2$ independent of $w_1$.  For every $v\in U$,
symmetry gives
\[
 (A_xv)\otimes w_2
   =A(u\otimes w_1,v\otimes w_2)
   =(A_{v\otimes w_2}u)\otimes w_1.
\]
The two indicated tensor subspaces have zero intersection, so
$A_xv=0$ for every $v$.  Thus $A_x=0$ for every pure tensor $x$,
and hence for every $x$.  No basis of $W$ invariant under any
larger Lie algebra has been assumed.
\end{proof}

\begin{lemma}[Nonzero prolongation produces a rank-one element]
\label{loc:rank-one-symbol}
Let $\mathfrak h\subset\End(T)$ be semisimple, faithful and
irreducible.  If $\mathfrak h^{(1)}\ne0$, then
$\mathfrak h$ contains an endomorphism of rank one.
\end{lemma}

\begin{proof}
Use homogeneous polynomial vector fields on the affine space $T$.
Let $\mathfrak p_{-1}=T$ be the translations,
$\mathfrak p_0=\mathfrak h$, and for $k\ge1$ set
\[
 \mathfrak p_k=
 \{A\in\Sym^{k+1}T^*\otimes T:
   A(x_1,\ldots,x_k,{-})\in\mathfrak h
                       \text{ for every }x_1,\ldots,x_k\in T\}.
\]
These spaces form a graded Lie algebra
$\mathfrak p=\bigoplus_{k\ge-1}\mathfrak p_k$ under the
vector-field bracket.  Indeed, differentiating a bracket gives the
sum of brackets with derivatives.  Induction on the total degree,
starting with closure of $\mathfrak h$ under brackets, proves
$[\mathfrak p_i,\mathfrak p_j]\subset\mathfrak p_{i+j}$.
The representation also gives
$[\mathfrak h,\mathfrak p_k]\subset\mathfrak p_k$ directly.

Suppose first that $\mathfrak p$ is finite-dimensional.  Its
solvable radical is graded, because it is preserved by the scaling
automorphisms of this grading.  From any nonzero homogeneous
polynomial vector field in that radical, successive brackets with
translations produce a nonzero constant field.  Irreducibility
under $\mathfrak h$ then puts all of $T$ in the radical.  Since
$\mathfrak p_1\ne0$, the bracket
$[T,\mathfrak p_1]$ is nonzero.  It lies in the intersection of
the radical with $\mathfrak h$, a solvable ideal of a semisimple
algebra, which is impossible.  Thus the radical is zero.

The grading derivation $D$, acting as $k\id$ on
$\mathfrak p_k$, must now be inner.  Recall a short proof of this
fact.  For any derivation $D$ of a semisimple algebra $\mathfrak a$,
nondegeneracy of the Killing form gives $z\in\mathfrak a$ such that
$E=D-\operatorname{ad}(z)$ satisfies
$\tr(E\operatorname{ad}(x))=0$ for all $x$.  Then
\[
 B(Ex,y)=\tr([E,\operatorname{ad}(x)]\operatorname{ad}(y))
        =\tr(E\operatorname{ad}([x,y]))=0,
\]
so $E=0$.  Apply this to $\mathfrak a=\mathfrak p$ and write
$D=\operatorname{ad}(z)$.  Comparing homogeneous degrees and
using the zero center shows $z\in\mathfrak p_0=\mathfrak h$.
But $D|_T=-\id_T$.  This would give a nonzero scalar in the
semisimple subalgebra $\mathfrak h\subset\mathfrak{sl}(T)$,
again impossible.  Hence $\mathfrak p$ is infinite-dimensional.

We make the symbol argument for this last alternative explicit.
Let $S=\Sym T$, and let
$\mathfrak h^\perp\subset T\otimes T^*$ be the annihilator of
$\mathfrak h\subset T^*\otimes T$.  Form the finitely generated
graded $S$-module
\[
 \mathcal M=(S\otimes T^*)/
                    S\mathfrak h^\perp,
\]
where the relations have degree one.  Perfect pairing with
homogeneous polynomial vector fields identifies
$\mathcal M_d^*=\mathfrak p_{d-1}$ for $d\ge0$:
annihilating a relation multiplied by $d-1$ linear factors is
exactly the displayed partial-evaluation condition.  Since
$\mathfrak p$ is infinite-dimensional, $\mathcal M$ is not a
finite-length module supported at the origin.  Its support thus
contains a nonzero point $\xi\in T^*=\Spec(S)(\C)$.

The fiber at $\xi$ is
\[
 T^*/\{(\xi\otimes\id)(r):r\in\mathfrak h^\perp\},
\]
and is nonzero.  Choose a nonzero vector $u\in T$ annihilating
its relation space.  Then
$\langle r,u\otimes\xi\rangle=0$ for every
$r\in\mathfrak h^\perp$.  Hence $u\otimes\xi\in\mathfrak h$.
Since both $u$ and $\xi$ are nonzero, this is the required
rank-one endomorphism.
\end{proof}

\begin{lemma}[The rank-one root calculation]
\label{loc:rank-one-classification}
A faithful irreducible representation $T$ of a complex semisimple
Lie algebra containing a rank-one endomorphism is the standard
representation, or its dual, of $\mathfrak{sl}(T)$, or the
standard representation of $\mathfrak{sp}(T)$.  In particular the
Lie algebra is simple.
\end{lemma}

\begin{proof}
We give the weight argument, including the exceptional short-root
possibilities.  Consider in $\PP(\mathfrak h)$ the closure of
the orbit of a rank-one operator.  It consists of rank-one
operators, since the rank-at-most-one locus is closed and the zero
operator does not occur in projective space.  A Borel subgroup has
a fixed point on this closure.  In this particular linear setting
one can see the fixed point directly: triangularize the solvable
Borel action, choose the smallest term of an invariant complete
flag whose projectivization meets the closure, and intersect there.
This projective intersection avoids the preceding hyperplane and
is therefore also affine.  It is finite; a connected group fixes
each of its points.  In the adjoint representation a Borel-fixed
line is the highest-root line of one simple ideal.  The highest
root is long.  We have therefore found a highest-long-root
operator of rank one in one simple ideal.

An irreducible representation of a direct sum of simple Lie
algebras is an external tensor product of their irreducible
representations.  The rank of this root operator on $T$ is its
rank on the indicated factor multiplied by the dimensions of all
other factors.  Rank one forces those other factors to have
dimension one.  They act trivially, and faithfulness eliminates
them.  Thus $\mathfrak h$ is simple.

All long roots are Weyl-conjugate.  For every long root $\alpha$,
the raising operator $e_\alpha$ has rank one.  Restriction to the
root $\mathfrak{sl}_2$ is a direct sum of irreducibles; a summand
of highest weight $m$ contributes rank $m$ to $e_\alpha$.
There is consequently exactly one two-dimensional standard
summand and otherwise only trivial summands.  The Cartan
subalgebra commuting with this root $\mathfrak{sl}_2$ acts on
the unique standard summand by a single character.  Thus the root
reflection $s_\alpha$ interchanges exactly two full Cartan
weights, each of multiplicity one and differing by $\alpha$,
and fixes every other full weight.  This assertion does not yet
exclude a zero weight of higher multiplicity.

Consider one irreducible component of the subsystem of long roots.
Its simple reflections act by distinct transpositions on the set
of weights.  Adjacent simple roots have a product of reflections
of order three, so their transpositions share exactly one vertex.
Orthogonal simple roots give commuting transpositions.  Two
different simple roots cannot give the same transposition, since
the difference of the transposed weights would be both roots up
to sign.  Connectedness of the diagram therefore gives a single
connected moving set of vertices.  The transpositions generate
the full symmetric group on this set.

Map the augmentation representation on these vertices to the
weight space by sending a vertex difference to the corresponding
weight difference.  The map is nonzero and equivariant.  The
augmentation representation of a symmetric group is irreducible,
so the map is injective.  Its image is exactly the span of the
root component: consecutive vertex differences are its simple
roots, and all differences lie in that span.  Conjugating a
simple-root transposition shows that every pair difference is a
root.  Conversely every root in this component arises this way,
by conjugacy of its root reflections.  The component is
therefore of type $A$, with the moving weights the vertices of
its standard simplex, up to their common centroid.

Distinct long-root components have disjoint moving sets.
Their transpositions commute, so an overlap could only be an
identical pair; the difference of that pair would then belong to
two orthogonal root spans, which is impossible.  The long roots
span the entire Cartan dual.  Hence the centroid of a moving
simplex, being fixed by its component and by every other
component, is zero.  A weight moved by no long-root component is
zero.

The ordinary classification of reduced irreducible root systems
\cite[Theorem 3.10]{Cap2016} now leaves very few cases.
For completeness, their long-root
subsystems are
\[
\begin{array}{c|c}
 \text{root system}&\text{long-root subsystem}\\ \hline
 A_\ell&A_\ell\\
 B_\ell\ (\ell\ge3)&D_\ell\\
 C_\ell&A_1^{\ell}\\
 D_\ell\ (\ell\ge4)&D_\ell\\
 E_6,E_7,E_8&E_6,E_7,E_8\\
 F_4&D_4\\
 G_2&A_2
\end{array}
\]
Here $B_2=C_2$ and the rank-one cases are $A_1$.
For the non-simply-laced rows, this table follows directly from
the coordinate roots: the long roots of $B_\ell$ and $F_4$ are
$\pm\epsilon_i\pm\epsilon_j$, those of $C_\ell$ are
$\pm2\epsilon_i$, and the long roots of $G_2$ form a regular
hexagon.  In the simply-laced cases every root is long.
Thus only $A$, $C$, $B_3$ (where $D_3=A_3$), and $G_2$
are possible.  The Weyl groups of $B_3$ and $G_2$ contain
$-\id$.  In those cases the nonzero weights would be respectively
a tetrahedron or a triangle with centroid zero.  Such a simplex
is not preserved by $-\id$; adjoining zero weights cannot make
it centrally symmetric.  These two cases are excluded.

For type $A_\ell$, the remaining simplex weights are the standard
weights or their negatives.  The highest weight is therefore
the first or last fundamental weight.  Highest-weight uniqueness
\cite[Theorem 4.3]{Cap2016}
identifies $T$ with the standard module or its dual.  For type
$C_\ell$, the long roots are $\pm2\epsilon_j$; each centered
two-vertex simplex is $\{\epsilon_j,-\epsilon_j\}$.  The
highest weight is $\epsilon_1$, so highest-weight uniqueness
gives the standard symplectic module.  In both cases this last
uniqueness statement excludes any additional zero-weight
multiplicity.  The root-system and highest-weight facts used here
are the usual finite-dimensional characteristic-zero theorems;
the rank-one and simplex deductions have been proved above.
\end{proof}

\begin{theorem}[Structure of the residual tensor group]
\label{loc:slsp}
For a maximal chart there is a decomposition
\begin{equation}
 V_{\mathrm{fib}}={\bf1}^{\oplus s}\oplus\bigoplus_{i=1}^r L_i,
 \qquad
 G=\prod_{i=1}^rG_i,
 \qquad
 G_i=\operatorname{SL}(L_i)\ \text{or}\
       \operatorname{Sp}(L_i),
 \label{loc:group-product}
\end{equation}
where each $G_i$ acts by its standard representation on $L_i$
and trivially on all other summands.  There are no repeated
nontrivial irreducible summands.  We place
$\operatorname{SL}_2=\operatorname{Sp}_2$ in the symplectic list.
The integer $s$ is the number of free rank-one direct summands
of $V$.
\end{theorem}

\begin{proof}
By Proposition~\ref{loc:connected-semisimple},
$\mathfrak g$ is semisimple.  The one-block algebra
$\mathfrak h_i$ of Lemma~\ref{loc:block-obstruction} is an
ideal of $\mathfrak g$, hence semisimple.  Split off its
complementary ideal.  On the irreducible dual of the $i$th
block the representation is $U\otimes W$, with
$\mathfrak h_i$ acting on $U$.  If $\dim W>1$,
Lemma~\ref{loc:auxiliary-factor} makes the first prolongation
zero, contrary to Lemma~\ref{loc:block-obstruction} for a
nontrivial block.  Therefore $\mathfrak h_i$ acts irreducibly
there.  It acts faithfully, since its elements act by zero on
every other block and the total representation is faithful.
Lemmas~\ref{loc:rank-one-symbol} and
\ref{loc:rank-one-classification} identify it with the full
special linear or symplectic algebra in its standard module.

The ideals so obtained for distinct blocks commute, since each
is zero on all the other blocks.  A remaining complementary
semisimple ideal would commute with these irreducible standard
actions.  By Schur's lemma it would act by scalars on each block;
a semisimple Lie algebra has no scalar character.  It also acts
trivially on the trivial blocks, so faithfulness eliminates it.
If two nontrivial blocks were isomorphic as $\mathfrak g$-modules,
an element acting trivially on one would act trivially on the
other; their one-block ideals would be zero.  This again
contradicts Lemma~\ref{loc:block-obstruction}.

It remains to check the group, rather than just its Lie algebra.
Take the simply connected covering group of the connected
semisimple group $G$.  Its factors are the indicated
$\operatorname{SL}$ and $\operatorname{Sp}$ groups.  Each
factor acts faithfully by its standard representation on its
own block, and trivially on the others.  Thus the kernel on the
full representation is trivial, including any proposed central
or diagonal central subgroup.  The covering map is therefore
an isomorphism and its image is the blockwise product in
\eqref{loc:group-product}.  Finally a trivial summand in the
trace quotient is an actual free summand by
Lemma~\ref{loc:split-lift}, and conversely.  This identifies $s$
as claimed.
\end{proof}

\begin{corollary}[A unique exterior constituent]
\label{loc:exterior-constituent}
In the notation of Theorem~\ref{loc:slsp}, suppose $r>0$ and put
$W=\bigoplus_{i=1}^rL_i$.  There are a simple $G$-module $E$ and an
integer $1\leq j_0\leq\dim W$ such that, for $0\leq b\leq\dim W$,
\[
 \dim_{\C}\Hom_G(E,\Lambda^bW)=
 \begin{cases}1,&b=j_0,\\0,&b\ne j_0.\end{cases}
\]
The multiplicity of $E$ in $\Lambda^{j_0}V_{\mathrm{fib}}$ is also
one.  The selection of $E$ is intrinsic to the standard factors and
is preserved by permutations of the factors.
\end{corollary}

\begin{proof}
For a standard $\operatorname{SL}_m$ block with $m>2$, select its
degree-one representation.  The exterior powers in degrees
$1,\ldots,m-1$ are irreducible with distinct highest weights, while
degrees $0$ and $m$ are trivial.  The selected representation therefore
occurs once and in degree one only.

For a standard $\operatorname{Sp}_{2m}$ block $L_i$, select
\[
 P_i^m=\ker\bigl(\Lambda^mL_i\longrightarrow\Lambda^{m-2}L_i\bigr),
\]
where the arrow is symplectic contraction and negative exterior powers
are zero.  The primitive decomposition
\[
 \Lambda^bL_i\simeq
 \bigoplus_{a=0}^{\lfloor\min(b,2m-b)/2\rfloor}
 P_i^{\min(b,2m-b)-2a}
\]
has irreducible, mutually distinct primitive terms $P_i^s$ of highest
weight $\varpi_s$ \cite[Theorem~17.5]{FultonHarris1991}.
Thus $P_i^m$ occurs once and only for $b=m$.  This includes
$\operatorname{SL}_2=\operatorname{Sp}_2$.

Take the external tensor product of the selected representations.
It is simple for the product group and occurs in $\Lambda^\bullet W$
once, in a single degree $j_0$.  In $\Lambda^{j_0}V_{\mathrm{fib}}$
it still has multiplicity one: a positive exterior degree from the
trivial summand would require its occurrence in a smaller degree of
$W$.  The factorwise rule is preserved by permutations.  In the
symplectic case the kernel defining $P_i^m$ is also unchanged when
the symplectic form is rescaled, so the selection has the asserted
intrinsic meaning.
\end{proof}
